% !TEX program = tectonic
% =============================================================================
%  AULA 12 — Segurança em Redes sem Fio
%  Curso: Redes sem Fio  |  Profa. Anelise Munaretto
%  WPA2/WPA3, 802.1X, KRACK, IoT security, 5G AKA.
%  Compilar: tectonic aula12.tex
% =============================================================================
\documentclass[aspectratio=169,11pt]{beamer}
% =============================================================================
%  PREÁMBULO COMÚN para as aulas de Redes sem Fio (Beamer + TikZ)
%  Incluido com % =============================================================================
%  PREÁMBULO COMÚN para as aulas de Redes sem Fio (Beamer + TikZ)
%  Incluido com % =============================================================================
%  PREÁMBULO COMÚN para as aulas de Redes sem Fio (Beamer + TikZ)
%  Incluido com \input{preamble.tex} en cada aula.
%  Paleta de cores e estilos são definidos aqui uma única vez.
% =============================================================================

\usepackage[T1]{fontenc}
\usepackage{amsmath,amssymb,mathtools}
\usepackage{graphicx}
\usepackage{tikz}
\usetikzlibrary{positioning,arrows.meta,calc,backgrounds,decorations.pathmorphing,fit,shadows,shapes.symbols,shapes.geometric}
\usepackage{pgfplots}
\pgfplotsset{compat=1.16}
\usepackage{tabularx,booktabs,multirow,array}
\usepackage[normalem]{ulem}
\usepackage{hyperref}

% ---------- Paleta de cores própria ----------
\definecolor{aneliseAzul}{RGB}{20,60,110}
\definecolor{aneliseVerde}{RGB}{30,140,70}
\definecolor{aneliseLaranja}{RGB}{215,120,20}
\definecolor{aneliseGris}{RGB}{90,90,90}

% ---------- Tema ----------
\usetheme{Madrid}
\usecolortheme{whale}
\setbeamercolor{structure}{fg=aneliseAzul}
\setbeamercolor{title}{fg=aneliseAzul}
\setbeamercolor{frametitle}{fg=aneliseAzul}
\setbeamercolor{block title}{fg=aneliseAzul,bg=aneliseGris!12!white}
\setbeamercolor{block body}{bg=aneliseGris!7!white}
\hypersetup{colorlinks=true,linkcolor=aneliseAzul,urlcolor=aneliseVerde}

% ---------- Comandos auxiliares ----------
\newcommand{\kurose}{Kurose \& Ross, \emph{Computer Networking: a Top-Down Approach}}
\newcommand{\forouzan}{Forouzan, \emph{Data Communications and Networking}}
\newcommand{\stallings}{Stallings, \emph{Wireless Communications \& Networks}}
\newcommand{\tanenbaum}{Tanenbaum, \emph{Computer Networks}}
\newcommand{\rfc}[1]{\href{https://www.rfc-editor.org/rfc/rfc#1.txt}{RFC~#1}}
% -----------------------------------------------------------------------------
%  Estilos TikZ comuns para desenhar redes (posicionamento relativo)
% -----------------------------------------------------------------------------
\tikzset{
  host/.style={draw,rounded corners=2mm,fill=aneliseAzul!15,inner sep=1.5mm,font=\scriptsize},
  rt/.style={draw,rounded corners=1mm,fill=aneliseLaranja!25,inner sep=1mm,font=\scriptsize},
  nube/.style={draw,cloud,aspect=2.2,fill=aneliseGris!15,inner sep=1.5mm,font=\scriptsize},
  el/.style={draw,rounded corners=1.2mm,fill=aneliseGris!18,inner sep=0.8mm,font=\scriptsize},
  enl/.style={thick},
  enlA/.style={thick,-Stealth},
} en cada aula.
%  Paleta de cores e estilos são definidos aqui uma única vez.
% =============================================================================

\usepackage[T1]{fontenc}
\usepackage{amsmath,amssymb,mathtools}
\usepackage{graphicx}
\usepackage{tikz}
\usetikzlibrary{positioning,arrows.meta,calc,backgrounds,decorations.pathmorphing,fit,shadows,shapes.symbols,shapes.geometric}
\usepackage{pgfplots}
\pgfplotsset{compat=1.16}
\usepackage{tabularx,booktabs,multirow,array}
\usepackage[normalem]{ulem}
\usepackage{hyperref}

% ---------- Paleta de cores própria ----------
\definecolor{aneliseAzul}{RGB}{20,60,110}
\definecolor{aneliseVerde}{RGB}{30,140,70}
\definecolor{aneliseLaranja}{RGB}{215,120,20}
\definecolor{aneliseGris}{RGB}{90,90,90}

% ---------- Tema ----------
\usetheme{Madrid}
\usecolortheme{whale}
\setbeamercolor{structure}{fg=aneliseAzul}
\setbeamercolor{title}{fg=aneliseAzul}
\setbeamercolor{frametitle}{fg=aneliseAzul}
\setbeamercolor{block title}{fg=aneliseAzul,bg=aneliseGris!12!white}
\setbeamercolor{block body}{bg=aneliseGris!7!white}
\hypersetup{colorlinks=true,linkcolor=aneliseAzul,urlcolor=aneliseVerde}

% ---------- Comandos auxiliares ----------
\newcommand{\kurose}{Kurose \& Ross, \emph{Computer Networking: a Top-Down Approach}}
\newcommand{\forouzan}{Forouzan, \emph{Data Communications and Networking}}
\newcommand{\stallings}{Stallings, \emph{Wireless Communications \& Networks}}
\newcommand{\tanenbaum}{Tanenbaum, \emph{Computer Networks}}
\newcommand{\rfc}[1]{\href{https://www.rfc-editor.org/rfc/rfc#1.txt}{RFC~#1}}
% -----------------------------------------------------------------------------
%  Estilos TikZ comuns para desenhar redes (posicionamento relativo)
% -----------------------------------------------------------------------------
\tikzset{
  host/.style={draw,rounded corners=2mm,fill=aneliseAzul!15,inner sep=1.5mm,font=\scriptsize},
  rt/.style={draw,rounded corners=1mm,fill=aneliseLaranja!25,inner sep=1mm,font=\scriptsize},
  nube/.style={draw,cloud,aspect=2.2,fill=aneliseGris!15,inner sep=1.5mm,font=\scriptsize},
  el/.style={draw,rounded corners=1.2mm,fill=aneliseGris!18,inner sep=0.8mm,font=\scriptsize},
  enl/.style={thick},
  enlA/.style={thick,-Stealth},
} en cada aula.
%  Paleta de cores e estilos são definidos aqui uma única vez.
% =============================================================================

\usepackage[T1]{fontenc}
\usepackage{amsmath,amssymb,mathtools}
\usepackage{graphicx}
\usepackage{tikz}
\usetikzlibrary{positioning,arrows.meta,calc,backgrounds,decorations.pathmorphing,fit,shadows,shapes.symbols,shapes.geometric}
\usepackage{pgfplots}
\pgfplotsset{compat=1.16}
\usepackage{tabularx,booktabs,multirow,array}
\usepackage[normalem]{ulem}
\usepackage{hyperref}

% ---------- Paleta de cores própria ----------
\definecolor{aneliseAzul}{RGB}{20,60,110}
\definecolor{aneliseVerde}{RGB}{30,140,70}
\definecolor{aneliseLaranja}{RGB}{215,120,20}
\definecolor{aneliseGris}{RGB}{90,90,90}

% ---------- Tema ----------
\usetheme{Madrid}
\usecolortheme{whale}
\setbeamercolor{structure}{fg=aneliseAzul}
\setbeamercolor{title}{fg=aneliseAzul}
\setbeamercolor{frametitle}{fg=aneliseAzul}
\setbeamercolor{block title}{fg=aneliseAzul,bg=aneliseGris!12!white}
\setbeamercolor{block body}{bg=aneliseGris!7!white}
\hypersetup{colorlinks=true,linkcolor=aneliseAzul,urlcolor=aneliseVerde}

% ---------- Comandos auxiliares ----------
\newcommand{\kurose}{Kurose \& Ross, \emph{Computer Networking: a Top-Down Approach}}
\newcommand{\forouzan}{Forouzan, \emph{Data Communications and Networking}}
\newcommand{\stallings}{Stallings, \emph{Wireless Communications \& Networks}}
\newcommand{\tanenbaum}{Tanenbaum, \emph{Computer Networks}}
\newcommand{\rfc}[1]{\href{https://www.rfc-editor.org/rfc/rfc#1.txt}{RFC~#1}}
% -----------------------------------------------------------------------------
%  Estilos TikZ comuns para desenhar redes (posicionamento relativo)
% -----------------------------------------------------------------------------
\tikzset{
  host/.style={draw,rounded corners=2mm,fill=aneliseAzul!15,inner sep=1.5mm,font=\scriptsize},
  rt/.style={draw,rounded corners=1mm,fill=aneliseLaranja!25,inner sep=1mm,font=\scriptsize},
  nube/.style={draw,cloud,aspect=2.2,fill=aneliseGris!15,inner sep=1.5mm,font=\scriptsize},
  el/.style={draw,rounded corners=1.2mm,fill=aneliseGris!18,inner sep=0.8mm,font=\scriptsize},
  enl/.style={thick},
  enlA/.style={thick,-Stealth},
}

\title[Redes sem Fio]{Redes sem Fio\\ Segurança em Redes sem Fio}
\subtitle{Aula 12}
\author[Anelise Munaretto]{Profa. Anelise Munaretto}
\institute[Curso]{Redes sem Fio --- Ciência da Computação}
\date{2026}

\begin{document}

\begin{frame}[plain]
  \titlepage
\end{frame}

\begin{frame}{Objetivos de aprendizagem}
  Ao final desta aula, você será capaz de:
  \begin{itemize}
    \item \textbf{explicar} o ataque KRACK e \textbf{comparar} WPA2-PSK com WPA3-SAE;
    \item \textbf{descrever} a evolução WEP $\rightarrow$ WPA $\rightarrow$ WPA2 $\rightarrow$ WPA3 (criptografia e autenticação);
    \item \textbf{analisar} a segurança de LoRaWAN (OTAA vs.\ ABP) e do 5G AKA (SUCI, home control);
    \item \textbf{estimar} o tempo de ataque por força bruta a um PSK;
    \item \textbf{identificar} ameaças (\emph{evil twin}, dicionário offline, \emph{deauth}) e as mitigações correspondentes.
  \end{itemize}
  \begin{alertblock}{Pré-requisitos}
    Aulas 2 (fundamentos de redes sem fio) e 8 (IoT e redes LPWAN).
  \end{alertblock}
\end{frame}

\section{Protocolos de Segurança Wi-Fi}

\begin{frame}{Evolução da segurança em Wi-Fi}
  \begin{center}
    \scriptsize
    \begin{tabular}{@{}llll@{}}
      \toprule
      \textbf{Padrão} & \textbf{Criptografia} & \textbf{Autenticação} & \textbf{Ano}\\
      \midrule
      WEP & RC4 (64/128 bits) & Chave compartilhada & 1999\\
      WPA & TKIP/RC4 & PSK / 802.1X & 2003\\
      WPA2 & AES-CCMP & PSK / 802.1X & 2004\\
      WPA3 & AES-GCMP-256 & SAE (Simult. Auth. of Equals) & 2018\\
      WPA3-Enterprise & AES-GCMP-128/256 & EAP-TLS, SuiteB & 2018\\
      \bottomrule
    \end{tabular}
  \end{center}
  \begin{block}{WPA3: principais novidades}
    \begin{itemize}
      \item \textbf{SAE} (Dragonfly handshake): protege contra ataques de dicionário offline
      \item \textbf{Perfect Forward Secrecy} (PFS)
      \item \textbf{PMF} (Protected Management Frames): obrigatório
      \item \textbf{Wi-Fi Easy Connect} (DPP): onboarding seguro
    \end{itemize}
  \end{block}
  \note{WPA3-SAE (Dragonfly) elimina o ataque de dicionário offline porque o segredo é usado num protocolo de troca de chaves baseado em curva elíptica; PMF obrigatório bloqueia deauth forjado.}
\end{frame}

\begin{frame}{Ataques conhecidos}
  \begin{columns}
    \begin{column}{0.5\textwidth}
      \textbf{KRACK (2017)}
      \begin{itemize}
        \item Key Reinstallation Attack (Mathy Vanhoef)
        \item ataque ao 4-way handshake do WPA2
        \item força reuso de nonce $\rightarrow$ quebra criptografia
        \item afeta todos os clientes WPA2
        \item corrigido em patches (CVE-2017-13077 a 13084)
      \end{itemize}
      \textbf{Evil Twin / Rogue AP}
      \begin{itemize}
        \item AP malicioso com mesmo SSID
        \item captura de credenciais (phishing)
        \item WPA3-Enterprise mitiga com certificado de servidor
      \end{itemize}
    \end{column}
    \begin{column}{0.5\textwidth}
      \textbf{Dicionário offline}
      \begin{itemize}
        \item WPA2-PSK: capturar 4-way handshake e tentar senhas
        \item ferramentas: aircrack-ng, hashcat
        \item WPA3 SAE torna essa abordagem inviável
      \end{itemize}
      \textbf{Deauth Attack}
      \begin{itemize}
        \item envio de quadros de desautenticação falsos
        \item PMF obrigatório no WPA3 evita
      \end{itemize}
    \end{column}
  \end{columns}
  \note{KRACK explora a reinstalação da chave de sessão no 4-way handshake, forçando reuso de nonce; no WPA2-PSK o atacante captura o handshake (via deauth falso) e testa senhas offline com aircrack-ng/hashcat.}
\end{frame}

\section{Segurança IoT/LPWAN}

\begin{frame}{Segurança em LoRaWAN e LPWAN}
  \begin{itemize}
    \item \textbf{LoRaWAN}: criptografia AES-128 (NwkSKey, AppSKey)
      \begin{itemize}
        \item OTAA (Over-The-Air Activation) vs.\ ABP (Activation By Personalization)
        \item OTAA mais seguro: chaves renovadas a cada join
        \item vulnerabilidade: payload pode ser descriptografado se AppSKey é exposta
      \end{itemize}
    \item \textbf{NB-IoT / LTE-M}: segurança herdada do LTE (AES, SIM-based)
      \begin{itemize}
        \item autenticação mútua UE-rede (AKA)
        \item criptografia NAS e AS
        \item chaves armazenadas no SIM
      \end{itemize}
  \end{itemize}
  \begin{alertblock}{Desafio IoT}
    Dispositivos com recursos limitados (CPU, memória) dificultam atualizações de segurança e rotação de chaves.
  \end{alertblock}
\end{frame}

\section{Segurança 5G}

\begin{frame}{Segurança 5G}
  \begin{columns}
    \begin{column}{0.5\textwidth}
      \textbf{5G AKA (Authentication and Key Agreement)}
      \begin{itemize}
        \item evolução do EPS-AKA (4G)
        \item SUPI (identidade permanente) cifrado como SUCI
        \item \emph{home control}: rede doméstica sempre envolvida
        \item chaves derivadas: K, CK, IK, K$_{\text{AMF}}$, etc.
      \end{itemize}
    \end{column}
    \begin{column}{0.5\textwidth}
      \textbf{Inovações de segurança 5G}
      \begin{itemize}
        \item \textbf{SEPP}: proteção entre operadoras (signaling)
        \item \textbf{Network Slicing}: isolamento entre slices
        \item \textbf{MEC security}: execução confiável na borda
        \item \textbf{NR encryption}: AES-128/256, Snow3G, ZUC
      \end{itemize}
    \end{column}
  \end{columns}
  \note{Diferente do 4G, o 5G protege a identidade: o SUPI nunca é enviado em claro e aparece como SUCI cifrado; o AKA garante home control com a rede doméstica sempre no loop de autenticação.}
\end{frame}

\section{Desafios}

\begin{frame}{Desafios de segurança em redes sem fio}
  \begin{itemize}
    \item \textbf{Quantum computing}: ameaça a criptografia atual (RSA, ECC)
    \item \textbf{Supply chain}: backdoors em hardware/chipsets
    \item \textbf{Privacidade}: rastreamento de dispositivos (MAC aleatório, IMSI catcher)
    \item \textbf{Atualizações}: dispositivos IoT sem mecanismo de patch
    \item \textbf{Regulamentação}: Lei Geral de Proteção de Dados (LGPD), ePrivacy
  \end{itemize}
\end{frame}

\begin{frame}{Exercício resolvido --- força bruta sobre um PSK}
  \textbf{Problema:} estimar o tempo para quebrar por força bruta um PSK de 8 caracteres com alfabeto de 62 símbolos (A--Z, a--z, 0--9).
  \begin{enumerate}
    \item \textbf{Espaço de chaves:}
      $$62^8 = 218\,340\,105\,584\,896 \approx 2{,}18 \times 10^{14}.$$
    \item \textbf{Tempo a $10^5$ tentativas/s:}
      $$\frac{2{,}18 \times 10^{14}}{10^5} = 2{,}18 \times 10^9\ \text{s} \approx 69{,}2\ \text{anos}.$$
    \item \textbf{Comparação com dicionário:} com $10^8$ senhas comuns a $10^5$/s, o ataque leva apenas $10^3$ s $\approx 17$ min --- por isso senhas fracas são o ponto fraco do WPA2-PSK.
  \end{enumerate}
  \begin{alertblock}{Por que o WPA3-SAE (Dragonfly) muda o jogo?}
    No SAE o segredo não é comparável offline: cada tentativa exige um \emph{round-trip} com o AP (sujeito a \emph{rate-limiting}), inviabilizando o ataque de dicionário offline.
  \end{alertblock}
\end{frame}

\begin{frame}{Resumo da aula}
  \begin{itemize}
    \item \textbf{Wi-Fi}: WEP (RC4) $\rightarrow$ WPA (TKIP) $\rightarrow$ WPA2 (AES-CCMP) $\rightarrow$ WPA3 (AES-GCMP, SAE).
    \item \textbf{WPA3-SAE}: resistente a dicionário offline, PFS e PMF obrigatório.
    \item \textbf{KRACK (2017)}: reinstalação de chave no 4-way handshake do WPA2; corrigido por patches (CVE-2017-13077--13084).
    \item \textbf{Ameaças}: \emph{evil twin}/Rogue AP, dicionário offline, \emph{deauth} falso.
    \item \textbf{LoRaWAN}: AES-128 com NwkSKey/AppSKey; OTAA renova chaves a cada \emph{join} (mais seguro que ABP).
    \item \textbf{5G AKA}: SUPI cifrado como SUCI, \emph{home control}, chaves derivadas (K, CK, IK, K$_{\text{AMF}}$); SEPP protege o signaling entre operadoras.
    \item \textbf{Quantitativo}: $62^8 \approx 2{,}18 \times 10^{14}$; a $10^5$ tentativas/s $\rightarrow \approx 69$ anos.
  \end{itemize}
\end{frame}

\begin{frame}{Autoavaliação}
  \begin{enumerate}
    \item O ataque KRACK explora:
      \begin{itemize}
        \item (a) a reinstalação da chave de sessão no 4-way handshake do WPA2
        \item (b) a ausência de criptografia no WEP
        \item (c) o roubo da senha no servidor RADIUS
      \end{itemize}
      \pause
      \textbf{Gabarito:} (a)

    \item No WPA3, o PSK é substituído pelo protocolo:
      \begin{itemize}
        \item (a) EAP-TLS
        \item (b) SAE (Simultaneous Authentication of Equals)
        \item (c) TKIP/RC4
      \end{itemize}
      \pause
      \textbf{Gabarito:} (b)

    \item Em LoRaWAN, o modo OTAA é mais seguro que o ABP porque:
      \begin{itemize}
        \item (a) usa AES-64 em vez de AES-128
        \item (b) dispensa o servidor de rede
        \item (c) renova as chaves de sessão a cada \emph{join}
      \end{itemize}
      \pause
      \textbf{Gabarito:} (c)

    \item No 5G AKA, a identidade permanente do usuário (SUPI) é transmitida:
      \begin{itemize}
        \item (a) em texto claro no rádio
        \item (b) cifrada como SUCI
        \item (c) apenas em redes LTE
      \end{itemize}
      \pause
      \textbf{Gabarito:} (b)
  \end{enumerate}
\end{frame}

\begin{frame}{Laboratório sugerido}
  \begin{block}{Auditoria Wi-Fi controlada e legal}
    \begin{itemize}
      \item Montar um AP de laboratório próprio (sem acesso à rede de produção) e capturar o tráfego com Wireshark em modo monitor.
      \item Demonstrar o \textbf{4-way handshake}: identificar as mensagens EAPOL 1--4 na captura.
      \item Com aircrack-ng/hashcat, testar a PSK em \emph{wordlist} pequena para observar o ataque de dicionário offline ao WPA2-PSK.
      \item Repetir com WPA3 (se o AP suportar) e verificar que a captura offline não permite o mesmo ataque (SAE).
    \end{itemize}
  \end{block}
  \begin{alertblock}{Ética e legalidade}
    Executar apenas em equipamentos próprios e com autorização; interferir em redes alheias viola a lei (Lei 12.737/2012 -- Lei Carolina Dieckmann).
  \end{alertblock}
\end{frame}

\section{Bibliografia}

\begin{frame}{Bibliografia}
  \begin{itemize}
    \item IEEE 802.11i-2004 (WPA2); Wi-Fi Alliance, WPA3 Specification v3.0
    \item M. Vanhoef, F. Piessens, ``Key Reinstallation Attacks: Forcing Nonce Reuse in WPA2'', ACM CCS 2017
    \item 3GPP TS 33.501, \emph{Security architecture and procedures for 5G system}
    \item LoRa Alliance, \emph{LoRaWAN 1.1 Specification: Security}
    \item NIST SP 800-38B, \emph{Recommendation for Block Cipher Modes of Operation}
  \end{itemize}
\end{frame}

\begin{frame}{Leitura recomendada}
  \begin{itemize}
    \item IEEE 802.11i-2004 (WPA2) e Wi-Fi Alliance, \emph{WPA3 Specification v3.0} --- os dois padrões que definem o que vimos em sala.
    \item M. Vanhoef, F. Piessens, ``Key Reinstallation Attacks: Forcing Nonce Reuse in WPA2'', \emph{ACM CCS 2017} --- o artigo original do KRACK.
    \item 3GPP TS 33.501, \emph{Security architecture and procedures for 5G system} --- especificação oficial do 5G AKA e do SUPI/SUCI.
    \item LoRa Alliance, \emph{LoRaWAN 1.1 Specification: Security} --- OTAA/ABP e derivação de chaves.
    \item NIST SP 800-38B --- modos de operação de cifras de bloco (contexto do AES-GCM do WPA3).
  \end{itemize}
  \begin{alertblock}{Dica de estudo}
    Refaça o cálculo do exercício com senhas de 10 e 12 caracteres para ver como o espaço de chaves cresce exponencialmente.
  \end{alertblock}
\end{frame}

\end{document}